\rSec3[func.wrap.ref]{Non-owning wrapper}

\rSec4[func.wrap.ref.general]{General}

\pnum
The header provides partial specializations of \tcode{function_ref} for each combination of the possible replacements of the placeholders \cv{} and \placeholder{noex} where:
\begin{itemize}
\item \cv{} is either const or empty, and
\item \placeholder{noex} is either \tcode{true} or \tcode{false}.
\end{itemize}

\rSec4[func.wrap.ref.class]{Class template \tcode{function_ref}}

\begin{codeblock}
namespace std {
  template<class R, class... ArgTypes>
  class function_ref<R(ArgTypes...) @\cv{}@ noexcept(@\placeholder{noex}@)> {
  public:
    // \ref{func.wrap.ref.ctor}, constructors and assignment operators
    template<class F> function_ref(F*) noexcept;
    template<class F> constexpr function_ref(F&&) noexcept;
    template<auto c, class F> constexpr function_ref(constant_wrapper<c, F>) noexcept;
    template<auto c, class F, class U>
      constexpr function_ref(constant_wrapper<c, F>, U&&) noexcept;
    template<auto c, class F, class T>
      constexpr function_ref(constant_wrapper<c, F>, @\cv{}@ T*) noexcept;

    constexpr function_ref(const function_ref&) noexcept = default;
    constexpr function_ref& operator=(const function_ref&) noexcept = default;
    template<class T> function_ref& operator=(T) = delete;

    // \ref{func.wrap.ref.inv}, invocation
    @\added{constexpr }@R operator()(ArgTypes...) const noexcept(@\placeholder{noex}@);

  private:
    template<class... T>
      static constexpr bool @\exposidnc{is-invocable-using}@ = @\seebelownc@;             // \expos

    R (*@\exposidnc{thunk-ptr}@)(@\exposidnc{BoundEntityType}@, ArgTypes&&...) noexcept(@\placeholdernc{noex}@);      // \expos
    @\exposidnc{BoundEntityType}@ @\exposidnc{bound-entity}@;                                       // \expos
  };

  // \ref{func.wrap.ref.deduct}, deduction guides
  template<class F>
    function_ref(F*) -> function_ref<F>;
  template<auto c, class F0>
    function_ref(constant_wrapper<c, F0>) -> function_ref<@\seebelow@>;
  template<auto c, class F, class T>
    function_ref(constant_wrapper<c, F>, T&&) -> function_ref<@\seebelow@>;
}
\end{codeblock}

\pnum
An object of class \tcode{function_ref<R(Args...) \cv{} noexcept(\placeholder{noex})>} stores a pointer to function \exposid{thunk-ptr} and an object \exposid{bound-entity}. The object \exposid{bound-entity} has an unspecified trivially copyable type \exposid{BoundEntityType}, that models \libconcept{copyable} and is capable of storing a pointer to object value or a pointer to function value. The type of \exposid{thunk-ptr} is \tcode{R(*)(\exposidnc{BoundEntityType}, Args\&\&...) noexcept(\placeholder{noex})}.

\pnum
Each specialization of \tcode{function_ref} is a trivially copyable type\iref{term.trivially.copyable.type} that models \libconcept{copyable}.

\pnum
Within subclause~\ref{func.wrap.ref}, \tcode{\placeholder{call-args}} is an argument pack with elements such that
\begin{codeblock}
decltype((@\placeholder{call-args}@))...
\end{codeblock}
denote \tcode{ArgTypes\&\&...} respectively.

\rSec4[func.wrap.ref.ctor]{Constructors and assignment operators}

\begin{itemdecl}
template<class... T>
  static constexpr bool @\exposid{is-invocable-using}@ = @\seebelow@;
\end{itemdecl}

\begin{itemdescr}
\pnum
If \placeholder{noex} is \tcode{true}, \tcode{\exposid{is-invocable-using}<T...>} is equal to:
\begin{codeblock}
is_nothrow_invocable_r_v<R, T..., ArgTypes...>
\end{codeblock}
Otherwise, \tcode{\exposid{is-invocable-using}<T...>} is equal to:
\begin{codeblock}
is_invocable_r_v<R, T..., ArgTypes...>
\end{codeblock}
\end{itemdescr}

\begin{itemdecl}
template<class F>
  static constexpr bool @\exposid{is-convertible-from-specialization}@ = @\seebelow@;
\end{itemdecl}

\begin{itemdescr}
\pnum
If \tcode{F} denotes a specialization \tcode{function_ref<R(Args...) \cvqual{cv2} noexcept(\placeholder{noex2})>} for some placeholders \cvqual{cv2} and \placeholder{noex2}, \tcode{\exposid{is-convertible-from-specialization}<F>} is equal to:
\begin{codeblock}
is_convertible_v<R(&)(Args...) noexcept(@\placeholder{noex2}@), R(&)(Args...) noexcept(@\placeholder{noex}@)> &&
is_convertible_v<int @\cv{}@&, int @\cvqual{cv2}@&>.
\end{codeblock}
Otherwise, \tcode{\exposid{is-convertible-from-specialization}<F>} is \tcode{false}.
\end{itemdescr}

\begin{itemdecl}
template<class F> function_ref(F* f) noexcept;
\end{itemdecl}

\begin{itemdescr}
\pnum
\constraints
\begin{itemize}
\item \tcode{is_function_v<F>} is \tcode{true}, and
\item \tcode{\exposid{is-invocable-using}<F>} is \tcode{true}.
\end{itemize}

\pnum
\expects
\tcode{f} is not a null pointer.

\pnum
\effects
Initializes \exposid{bound-entity} with \tcode{f}, and \exposid{thunk-ptr} with the address of a function \tcode{\placeholder{thunk}} such that \tcode{\placeholder{thunk}(\exposid{bound-entity}, \placeholder{call-args}...)} is expression-equivalent\iref{defns.expression.equivalent} to \tcode{invoke_r<R>(f, \placeholder{call-args}...)}.
\end{itemdescr}

\begin{itemdecl}
template<class F> constexpr function_ref(F&& f) noexcept;
\end{itemdecl}

\begin{itemdescr}
\pnum
Let \tcode{T} be \tcode{remove_reference_t<F>}.

\pnum
\constraints
\begin{itemize}
\item \tcode{remove_cvref_t<F>} is not the same type as \tcode{function_ref},
\item \tcode{is_member_pointer_v<T>} is \tcode{false}, and
\item \tcode{\exposid{is-invocable-using}<\cv{} T\&>} is \tcode{true}.
\end{itemize}

\pnum
\effects
If \tcode{\exposid{is-convertible-from-specialization}<remove_cv_t<T>>} is \tcode{false}, initializes \exposid{bound-entity} with \tcode{addressof(f)}, and \exposid{thunk-ptr} with the address of a function \tcode{\placeholder{thunk}} such that \tcode{\placeholder{thunk}(\exposid{bound-entity}, \placeholder{call-args}...)} is expression-equivalent\iref{defns.expression.equivalent} to \tcode{invoke_r<R>(static_cast<\cv{} T\&>(f), \placeholder{call-args}...)}. Otherwise, initializes \exposid{bound-entity} with the value of \tcode{f.\exposid{bound-entity}} and \exposid{thunk-ptr} with the value of \tcode{f.\exposid{thunk-ptr}}.

\pnum
\remarks
If \tcode{remove_cvref_t<F>} is a specialization of \tcode{function_ref}, an implementation may initialize \exposid{bound-entity} with the value of \tcode{f.\exposid{bound-entity}} and \exposid{thunk-ptr} with the value of \tcode{f.\exposid{thunk-ptr}}.
\begin{example}
\begin{codeblock}
void f1(std::string);
void f2(std::string);

function_ref<void(std::string)> r1(&f1);
function_ref<void(std::string&&)> r2(r1);
r2("");                 // \tcode{f1} is invoked
r1 = &f2;
r2("");                 // it is unspecified if \tcode{f1} or \tcode{f2} is invoked
\end{codeblock}
\end{example}
\end{itemdescr}

\begin{itemdecl}
template<auto c, class F> constexpr function_ref(constant_wrapper<c, F> f) noexcept;
\end{itemdecl}

\begin{itemdescr}
\pnum
\constraints
\tcode{\exposid{is-invocable-using}<const F\&>} is \tcode{true}.

\pnum
\mandates
\begin{itemize}
\item If \tcode{is_pointer_v<F> || is_member_pointer_v<F>} is \tcode{true}, then \tcode{f.value != nullptr} is \tcode{true}, and
\item if \tcode{ArgTypes} is not an empty pack and all types in \tcode{remove_cvref_t<ArgTypes>...} satisfy \exposconcept{constexpr-param} then \tcode{constant_wrapper<\placeholdernc{INVOKE}(f.value, remove_cvref_t<ArgTypes>::\newline value...)>} is not a valid type.
\end{itemize}

\pnum
\effects
Initializes \exposid{bound-entity} with a pointer to an unspecified object or null pointer value, and \exposid{thunk-ptr} with the address of a function \tcode{\placeholder{thunk}} such that \tcode{\placeholder{thunk}(\exposid{bound-entity}, \placeholder{call-args}...)} is expression-equivalent\iref{defns.expression.equivalent} to \tcode{invoke_r<R>(f.value, \placeholder{call-args}...)}.
\end{itemdescr}

\begin{itemdecl}
template<auto c, class F, class U>
  constexpr function_ref(constant_wrapper<c, F> f, U&& obj) noexcept;
\end{itemdecl}

\begin{itemdescr}
\pnum
Let \tcode{T} be \tcode{remove_reference_t<U>}.

\pnum
\constraints
\begin{itemize}
\item \tcode{is_rvalue_reference_v<U\&\&>} is \tcode{false}, and
\item \tcode{\exposid{is-invocable-using}<const F\&, \cv{} T\&>} is \tcode{true}.
\end{itemize}

\pnum
\mandates
If \tcode{is_pointer_v<F> || is_member_pointer_v<F>} is \tcode{true}, then \tcode{f.value != nullptr} is \tcode{true}.

\pnum
\effects
Initializes \exposid{bound-entity} with \tcode{addressof(obj)}, and \exposid{thunk-ptr} with the address of a function \tcode{\placeholder{thunk}} such that \tcode{\placeholder{thunk}(\exposid{bound-entity}, \placeholder{call-args}...)} is expression-equivalent\iref{defns.expression.equivalent} to \tcode{invoke_r<R>(f.value, static_cast<\cv{} T\&>(obj), \placeholder{call-args}...)}.
\end{itemdescr}

\begin{itemdecl}
template<auto c, class F, class T>
  constexpr function_ref(constant_wrapper<c, F> f, @\cv{}@ T* obj) noexcept;
\end{itemdecl}

\begin{itemdescr}
\pnum
\constraints
\tcode{\exposid{is-invocable-using}<const F\&, \cv{} T*>} is \tcode{true}.

\pnum
\mandates
If \tcode{is_pointer_v<F> || is_member_pointer_v<F>} is \tcode{true}, then \tcode{f.value != nullptr} is \tcode{true}.

\pnum
\expects
If \tcode{is_member_pointer_v<F>} is \tcode{true}, \tcode{obj} is not a null pointer.

\pnum
\effects
Initializes \exposid{bound-entity} with \tcode{obj}, and \exposid{thunk-ptr} with the address of a function \tcode{\placeholder{thunk}} such that \tcode{\placeholder{thunk}(\exposid{bound-entity}, \placeholder{call-args}...)} is expression-equivalent\iref{defns.expression.equivalent} to \tcode{invoke_r<R>(f.value, obj, \placeholder{call-args}...)}.
\end{itemdescr}

\begin{itemdecl}
template<class T> function_ref& operator=(T) = delete;
\end{itemdecl}

\begin{itemdescr}
\pnum
\constraints
\begin{itemize}
\item \tcode{\exposid{is-convertible-from-specialization}<T>} is \tcode{false},
\item \tcode{is_pointer_v<T>} is \tcode{false}, and
\item \tcode{T} is not a specialization of \tcode{constant_wrapper}.
\end{itemize}
\end{itemdescr}

\rSec4[func.wrap.ref.inv]{Invocation}

\begin{itemdecl}
@\added{constexpr }@R operator()(ArgTypes... args) const noexcept(@\placeholder{noex}@);
\end{itemdecl}

\begin{itemdescr}
\pnum
\effects
Equivalent to: \tcode{return \exposid{thunk-ptr}(\exposid{bound-entity}, std::forward<ArgTypes>(args)...);}
\end{itemdescr}

\rSec4[func.wrap.ref.deduct]{Deduction guides}

\begin{itemdecl}
template<class F>
  function_ref(F*) -> function_ref<F>;
\end{itemdecl}

\begin{itemdescr}
\pnum
\constraints
\tcode{is_function_v<F>} is \tcode{true}.
\end{itemdescr}

\begin{itemdecl}
template<auto c, class F0>
  function_ref(constant_wrapper<c, F0>) -> function_ref<@\seebelow@>;
\end{itemdecl}

\begin{itemdescr}
\pnum
Let \tcode{F} be \tcode{remove_pointer_t<F0>}.

\pnum
\constraints
\tcode{is_function_v<F>} is \tcode{true}.

\pnum
\remarks
The deduced type is \tcode{function_ref<F>}.
\end{itemdescr}

\begin{itemdecl}
template<auto c, class F, class T>
  function_ref(constant_wrapper<c, F>, T&&) -> function_ref<@\seebelow@>;
\end{itemdecl}

\begin{itemdescr}
\pnum
\constraints
\begin{itemize}
\item \tcode{F} is of the form \tcode{R(G::*)(A...) \cv{} \opt{\&} noexcept(E)} for a type \tcode{G}, or
\item \tcode{F} is of the form \tcode{M G::*} for a type \tcode{G} and an object type \tcode{M}, in which case let \tcode{R} be \tcode{invoke_result_t<F, T\&>}, \tcode{A...} be an empty pack, and \tcode{E} be \tcode{true}, or
\item \tcode{F} is of the form \tcode{R(*)(G, A...) noexcept(E)} for a type \tcode{G}.
\end{itemize}

\pnum
\remarks
The deduced type is \tcode{function_ref<R(A...) noexcept(E)>}.
\end{itemdescr}
